% Part: first-order-logic
% Chapter: beyond
% Section: other-logics

\documentclass[../../../include/open-logic-section]{subfiles}

\begin{document}

% SEGMENT OLP-0181-S01

\olfileid{fol}{byd}{oth}

\olsection{பிற தருக்கங்கள்}

இதற்குள் நீங்கள் உணர்ந்திருக்கக்கூடியபடி, புதிய தருக்கம் ஒன்றை
வடிவமைப்பது கடினமல்ல. நீங்களும்கூட உங்களுக்கென ஒரு தொடரியலை
உருவாக்கி, ஓர் உய்த்தறி அமைப்பைப் புனைந்து, அதனுடன் பொருந்தும்
பொருண்மையியலை அமைக்கலாம். அப்பொருண்மையியலுக்கு !!{derivation}
அமைப்பு நிறைவானதாக இருக்க வேண்டுமெனில் நீங்கள் சற்றுத் திறமையாகச்
செயல்பட வேண்டியிருக்கலாம்; உங்கள் தருக்கம் உண்மையிலேயே சுவையானது என்று
உலகெங்கும் உள்ளவர்களை நம்பச் செய்யவும் முயற்சி தேவைப்படலாம். ஆனால்
அதற்குப் பலனாக, உங்கள் தருக்கத்தின் கணித மற்றும் கணிப்பியல் பண்புகளை
ஆராய்ந்து பல மணிநேரம் நல்ல, தூய மகிழ்ச்சியை அனுபவிக்கலாம்.

அண்மைப் பத்தாண்டுகளில் முறைசார் தருக்கங்கள் உண்மையிலேயே வெடித்துப்
பெருகியுள்ளன. தெளிவற்ற பண்புகளைப் பற்றிய ஆய்வுமுறையை மாதிரியாக்க
தெளிவிலித் தருக்கம் வடிவமைக்கப்பட்டுள்ளது. ஐயப்பாட்டைப் பற்றிய
ஆய்வுமுறையை மாதிரியாக்க நிகழ்தகவுத் தருக்கம் வடிவமைக்கப்பட்டுள்ளது.
புதுத் தகவலை எதிர்கொள்ளும்போது பின்னர் மாற்றிவிடக்கூடிய
``நியாயமான'' உய்த்தறிதல்களாகிய முறியக்கூடிய ஆய்வுமுறைகளை மாதிரியாக்க,
இயல்புநிலைத் தருக்கங்களும் ஓரியக்கமற்ற தருக்கங்களும்
வடிவமைக்கப்பட்டுள்ளன. அறிவைப் பற்றிய ஆய்வுமுறையை மாதிரியாக்க
அறிவுநிலைத் தருக்கங்கள் உள்ளன; காரண உறவுகளைப் பற்றிய ஆய்வுமுறையை
மாதிரியாக்க காரணத் தருக்கங்கள் உள்ளன; அறம் மற்றும் ஒழுக்கக்
கடப்பாடுகளைப் பற்றிய ஆய்வுமுறையை மாதிரியாக்க வடிவமைக்கப்பட்ட
``கடப்பாட்டியல்'' தருக்கங்களும்கூட உள்ளன. இவ்வமைப்புகளை அறிமுகப்படுத்தும்
முதன்மை உந்துதல் மெய்யியல்சார்ந்ததா, கணிதம்சார்ந்ததா,
கணிப்பியல்சார்ந்ததா என்பதைப் பொறுத்து, கணிதத் தருக்கம், மெய்யியல்
தருக்கம், செயற்கை நுண்ணறிவு, அறிவாற்றல் அறிவியல் அல்லது வேறொரு
துறையின் கீழ் இவ்வுயிரினங்கள் படிக்கப்படுவதை நீங்கள் காணலாம்.

பட்டியல் நீண்டுகொண்டே செல்கிறது; சாத்தியப்பாடுகள் முடிவற்றவையாகத்
தோன்றுகின்றன. மனித ஆய்வுமுறை அனைத்தையும் கணக்கீடாகக் குறைக்க வேண்டும்
என்ற லைப்னிட்சின் கனவை நாம் ஒருவேளை எப்போதும் அடையாமல் போகலாம்---ஆனால்
முயற்சி செய்வதிலிருந்து அது நம்மைத் தடுக்க முடியாது.

\end{document}
